\section{The affine part and endpoint potentials}\label{sec:affine}

The state cochain now separates into a spatially affine term and two
local endpoint terms. Short comparison experiments control the endpoint
terms uniformly; the geometry of a long connector enters only a
controlled deterministic remainder.

Fix a finite normalized menu.  Actual placements may be its translates
on the $h$-grid; auxiliary translated and rotated copies are also allowed.
At fixed $h$, pass to a further subsequence on which the admissions
stabilize.  If any type is admitted, choose one reference type and its
color flip.  When the admitted menu is empty all comparisons at qualified
endpoints are vacuous; we may set $a_h^V=1$ and choose any bounded $a_h^Q$.

\subsection{Fixed stubs and straight means}

For each type fix two disjoint \emph{stubs} outside its obstacle, taking
the departure port to an eastward horizontal strip and the arrival port
to a westward horizontal strip.  Their exit sections lie at the midpoints
of opposite sides of a sufficiently large square, with transverse
coordinate $v$ increasing upward.  To construct them, use a Jordan
annulus between the obstacle and the square, match the two attachment
intervals in cyclic order, and draw disjoint product continuations.
The transverse order is along the positive boundary direction at the
east exit and against it at the west exit, precisely as required by the
departure and arrival ports.  Continue the strips straight outward.

Fix the local portions of the stubs and their positive margins in all
regions where chart changes may be supported before placing other
plugs.  A later enlargement of the exit square uses only straight
continuation.  Rotate a type together with its stubs whenever needed.
These choices meet the local-template hypotheses of
Lemma~\ref{lemma:chart-stability}.

\begin{proposition}[Affine straight means]\label{prop:affine}
For $X=Q,V$ there is a scalar $a_h^X$, uniformly bounded in $h$, such
that every fixed straight transport in the countable family used below
has mean
\begin{equation}\label{affine:means}
 \EE_\mu i_e^\eps=a_h^X(z_t-z_s)+o(1).
\end{equation}
The same scalar applies to translations, common rotations and color
flips of the reference type.  All choices can be made simultaneously
for the two fields and countably many $h$.
\end{proposition}

\begin{proof}
Join identical unrotated reference copies on a horizontal line, using
their fixed stubs and a straight gap.  Their center separation is
$\ell h$, where $\ell>D$ for a fixed sufficiently large $D$.  Write
$A_\eps^X(\ell)=\EE_\mu i_e^\eps$.  A remote complementary connector
through the unused stubs closes both this transport of length
$(\ell+r)h$ and the two transports obtained by inserting a middle
reference copy.  Lemma~\ref{lemma:common-closure} gives
\[
 A_\eps^X(\ell+r)=A_\eps^X(\ell)+A_\eps^X(r)+o(1),
       \qquad \ell,r>D.
\]
Translation and rotation of the entire experiment, including the root
ordering, preserve its law and rotate both fields as vectors.  If a
conditional-average or gate tie-breaking convention was chosen
separately at the translated placement, its transported version is an
alternative defining completion.  Gate consistency with common gates
in the shared states and Proposition~\ref{prop:state-cochain} compare
the two choices to $o_{L^1}(1)$.  Thus covariance to this accuracy is
sufficient.  Color exchange preserves the law as well.

The ratios $A_\eps^X(\ell)/h$ are bounded uniformly for
$D+1\le\ell\le D+4$, in the limit superior: these transports have
bounded normalized size and uniformly buffered short completions.
For each fixed larger length there is likewise a finite bound, not
necessarily uniform in that length.

Only countably many lengths are required: the standard routes below
are fixed by grid terminal pairs, finitely many types and integer twists.
Include the positive integers and take the rational span of these
lengths, a countable additive subgroup $\mathcal L\subset\R$.  Extract
limits of $A_\eps^X(\ell)/h$ for all $\ell\in\mathcal L$ with $\ell>D$.
Their function $A$ is additive whenever both arguments exceed $D$.
Extend it to $x\in\mathcal L$ by
\[
 \widetilde A(x)=A(u+x)-A(u),\qquad u,u+x>D.
\]
The value is independent of a sufficiently large $u$: comparing two
choices $u,v$ uses
$A(u+x)+A(v)=A(u+v+x)=A(v+x)+A(u)$.
Choosing all offsets large similarly proves additivity on $\mathcal L$.
For small $x$, take a fixed rational $u$ in the interior of
$[D+1,D+4]$; the interval bound gives a common bound on
$\widetilde A(x)$.  Additivity then gives continuity at zero: if
$|x|<\delta/N$, a bound for $\widetilde A(Nx)$ gives the same bound
divided by $N$ for $\widetilde A(x)$.  Since $\mathcal L$ contains the
rationals, continuity implies $\widetilde A(x)=a_h^X x$.  The interval
bound also bounds $|a_h^X|$ independently of $h$.  Covariance yields
\eqref{affine:means} in every direction.  Countable diagonal extraction
makes these choices simultaneously as asserted.
\end{proof}

\subsection{Normalization of the tangent field}

\begin{lemma}\label{lemma:tangent-normalization}
The scalar of the tangent field is
\begin{equation}\label{affine:tangent}
 a_h^V=1.
\end{equation}
\end{lemma}

\begin{proof}
We give the mask argument also for curvilinear end charts.  In an
unperturbed straight experiment choose a deterministic polygonal Jordan
quadrilateral $R_e$ approximating the sector from section $0$ to section
$0$ with transverse interval $[-2.7,1.5]$.  Its boundary is close in the
ordered parametrization.  The transverse end cuts remain in the first
gate balls, and its top and bottom sides lie beyond the possible rim
layer $[-2.35,.8]$ on its left and right, respectively.  Permit a small
longitudinal enlargement only inside the cutoff balls.  By compact
margins, outside slightly smaller cutoff balls the mask contains exactly
the relevant sector side of every confined rim, with room, and excludes
the complementary side.

Here the required polygonal approximation follows directly from Jordan
simplicity.  Trace a polygonal closed walk uniformly close to the
boundary in order, split any crossings in general position and retain
a simple subcycle with positive winding about a fixed inner point.
Dropped parameter gaps have endpoints tending together, and therefore
are short by injectivity of the limiting Jordan parametrization.  They
cannot be almost a whole turn: the retained cycle would then be local
and have zero winding about the inner point.  This preserves the ordered
approximation and the four chosen sides.  Fix one such quadrilateral in
a sufficiently long template.  For larger $\ell$, stretch only its
straight middle region horizontally, keeping both end regions rigid.
The margins persist for each fixed $\ell$, and the total length of the
two end cuts is $O(h)$ uniformly in $\ell$.

On good geometry, polygonal intersection with the left side of $P$
gives
\[
 \int_{P\cap R_e}\dd z=\ell h+O(h).
\]
Indeed this integral is minus the integral over the portion of
$\partial R_e$ on that left side.  The entire top side contributes with
reversed orientation, the bottom side does not contribute, and the end
cuts cost $O(h)$.  Their endpoint locations differ from the plug centers
by $O(h)$.  A fixed cut has probability zero of an accidental exact
overlap with a random primal edge or vertex degeneracy.

We next compare this geometric integral with the barycenter mask
$V_P(\ind_{R_e})$.  On fine mesh a discrepancy can involve only rim
edges whose tagged triangles meet a cut segment.  Put $s=\eps^\alpha$
with $\alpha<1$ sufficiently close to one and cover these finitely many
segments by $O_e(s^{-1})$ boxes of side $s$.  In any one enlarged box the
$L^2$ norm of the total candidate edge length is $O(s^2/\eps)$, by the
scaled local triangle moments.  The presence of a candidate rim edge
implies a monochromatic arm from scale $s$ to a fixed positive distance
for this pattern, since the rim also visits a different plug.  Its
ordinary probability is $O_e(s^c)$, apart from a negligible mesh event.
Cauchy--Schwarz and summation thus bound the expected discrepancy by
\[
 C_e s^{-1}\frac{s^2}{\eps}s^{c/2}
       =C_e\eps^{\alpha(1+c/2)-1}\longrightarrow0
\]
when $\alpha>1/(1+c/2)$.  Bad mesh events are discarded using the same
local moments.  The bounded density of this fixed experiment transfers
the assertion to its law.  Alternatively, convergence in probability
and the Bessel $L^2$ bound suffice for the corresponding mean limit.

Outside the two cutoff balls, $R_e$ gives exactly $\widehat j_e$.
Its difference from $\widehat m_s+\widehat j_e$ consists of deterministic
masks inside the endpoint cutoff balls.  Couple at either endpoint state
to a fixed short completion.  The inner second-gate paths and their
triangles agree, so these correction means are $O(h)+o(1)$ by the short
second-moment estimate.  Their constant is independent of $\ell$:
the end masks were fixed before the middle strip was stretched.  By
Proposition~\ref{prop:state-cochain} and \eqref{affine:means},
\[
 |a_h^V\ell-\ell|\le C
\]
for arbitrarily large positive integers $\ell$.  This proves
\eqref{affine:tangent}.
\end{proof}

\subsection{The two endpoint twists}

We next compare an arbitrary fixed connector to a controlled standard
representative.  Its class is understood with the end disks, port
points and charts fixed.

\begin{lemma}[Classification by endpoint twists]\label{lemma:two-twists}
After normalizing the two disjoint Jordan obstacles to round holes,
every exterior simple centerline can be replaced, by common-closure
comparisons, by a smooth proper arc radial at its ends.  Such arcs with
fixed endpoints are classified, for the comparisons needed here, by two
integer endpoint twists.  Relative to a standard route with bounded
endpoint turns, their total size $k$ is bounded by a constant plus a
constant times the sum of the absolute net argument changes about the
two centers along the exterior arc.
\end{lemma}

\begin{proof}
The smoothing and full-port collars were established after
Lemma~\ref{lemma:band-collar}.  Identify each hole boundary with the
circle of directions at its center.  Fix a radial identification of the
plane with the two punctures blown up and the exterior of the circular
holes: in disjoint neighborhoods of the centers, send radius zero to the
hole radius by an increasing radial coordinate, agreeing with the
identity outside those neighborhoods.  Apply its inverse to each
radial-ended proper arc and fill in the punctures.  The result is a
smooth embedded compact interval, radial near its ends, and the fixed
identification recovers exactly the original proper arc after blowing
up.  Any two such smooth intervals, with ordered endpoints
matched, are carried one to the other by an orientation-preserving plane
diffeomorphism. To obtain this smooth extension, first map smooth
strip-and-cap neighborhoods of the two intervals by an
orientation-preserving diffeomorphism, preserving their ordered axes.
Choose smooth disk boundaries inside these neighborhoods. The prescribed
Jordan--Schoenflies theorem extends this map across their complementary
surfaces, retaining its smooth values on collars of those boundaries.
Hatcher's smoothing theorem for surface homeomorphisms, in its form
relative to a boundary neighborhood
\cite[Theorem~B and the following refinement, p.~2]{Hatcher2025},
replaces the exterior homeomorphism by a diffeomorphism without changing
these collars. Gluing with the retained neighborhood map gives the
required plane diffeomorphism and preserves the two intervals and their
endpoints exactly.

An orientation-preserving plane diffeomorphism has a locally smooth
deformation to the identity.  After subtracting its value at a base
point, the maps $x\mapsto F(tx)/t$ tend on compact sets to its positive
differential as $t\downarrow0$; restore the translation and join the
positive linear map to the identity.  During the resulting deformation
of the interval, postcompose by the unique direct similarity restoring
the two current ordered endpoint positions to the two fixed punctures.
This correction is the identity at both ends of the deformation.
Thus the punctures stay fixed throughout.

Blowing up their directions gives a deformation of proper arcs whose
ends move on the two hole circles.  Lift both moving endpoint angles to
real-valued continuous functions.  Compensate each angle in a disjoint
boundary collar, tapering the compensating rotation to zero away from
that boundary.  The deformation now fixes the port points.  At its end
the two compensations are integral full twists, since the original and
final port points agree.  These are the only possible differences
between the fixed-end deformation classes.  The compensated arcs are
transverse to the hole circles.  Compactness of the parameter interval
provides uniformly small product collars in which all these transverse
germs can be straightened to radial germs, continuously in the
parameter.  Take the straightening correction to be the identity at
both terminal times, where the germs are already radial.  The resulting
deformation therefore has the full fixed-end collar form required for
the common-closure comparisons.

They are detected separately by net arguments.  For the starting
puncture, lift the argument along the deforming arc.  The difference of
its final and initial net argument changes is minus the net turn of the
starting direction loop: the direction at the other puncture has zero
net argument change about the starting puncture.  For the arriving
puncture the analogous identity has the opposite sign.  Hence equal
net argument changes at both centers make both compensations trivial
at the end, giving an endpoint-fixed deformation of the kind used in
Lemma~\ref{lemma:common-closure}.  Conversely integer full turns at the
ends realize the corresponding changes; the explicit disjoint
representatives are constructed below.

Choose the signs of the two integers so they record the additional
argument turns in the orientation from $s$ to $t$.  Their absolute sum
is bounded by the displayed net changes divided by $2\pi$, plus the
bounded turns of the reference route.  Normalizing the Jordan disks
does not change winding differences between two arcs with the same
ends: these differences are integer winding numbers of their joined
closed curves about the respective centers.  For allowed perturbed
charts, pull the route back first to the nominal disks.  The arc stays
outside the fixed central balls, and the small supported change has
only a bounded effect on the endpoint argument choices.  This gives
the same bound for $k$ up to a fixed additive constant.
\end{proof}

\subsection{Quantitative standard representatives}

For each terminal type choose a fixed short start insertion: go east
along its departure stub to a horizontally aligned reference copy at a
fixed large distance.  Choose a fixed short end insertion from a
reference copy west of the arriving terminal into its arrival stub.
Their distances are large enough for the full local structures and the
supports of chart changes; all these lengths are nevertheless fixed in
normalized units.  They have uniformly bounded short standalone
completions and hence the $C h$ cochain bounds.  The first and last
insertions will remain fixed when the long route changes.

\begin{lemma}[Standard spirals and short corner templates]
\label{lemma:standard-spirals}
Fix nominal terminal placements, their types and two twist integers
$n_s,n_t$, and put $k=|n_s|+|n_t|$.  There is a constant $C$ such that,
if
\[
 |z_t-z_s|>Ch(1+k),
\]
the class has a buffered representative beginning and ending with the
fixed insertions above.  Between them it consists of straight transports
and at most $C(1+k)$ short pieces.  The short pieces belong, up to rigid
motion, to uniformly buffered compact families of bounded normalized
size, and satisfy
\[
 \|i_j^\eps\|_{L^2}\le C_{\rm short}h
\]
for small $\eps$, with a constant independent of the terminal separation,
the twists and the particular long connector.

A common complementary connector closes both the inserted path and its
unsplit representative.  All auxiliary placements and straight lengths
can be fixed from the nominal terminal data and integers before making
allowed chart changes.  In the applications they form a countable list.
\end{lemma}


The three features of this representative have different roles in the
cochain comparison. Its fixed first and last insertions determine the
endpoint state functions. Its straight pieces, however long, have the
exact limiting means in Proposition~\ref{prop:affine}. Only the
$O(1+k)$ short pieces contribute an $O(h(1+k))$ deterministic error.
Figure~\ref{fig:standard-representative} shows how the used route and a
common complement can turn around the terminals without meeting.
We first derive the endpoint formula from these properties; the full
geometric construction follows after that proof.

\begin{figure}[tb]
\centering
\begin{tikzpicture}[x=1.45cm,y=1.45cm,>=Stealth,font=\small]
 \draw[gray!65,densely dotted] (0,0) circle[radius=1.46];
 \draw[gray!65,densely dotted] (4.4,0) circle[radius=1.46];
 \filldraw[fill=gray!15,draw=black] (0,0) circle[radius=.22];
 \filldraw[fill=gray!15,draw=black] (4.4,0) circle[radius=.22];
 \node at (0,0) {$s$};
 \node at (4.4,0) {$t$};
 \draw[very thick] (.22,0)--(.65,0);
 \draw[very thick,domain=0:360,samples=100,variable=\u]
       plot ({(.65+.70*\u/360)*cos(\u)},
             {(.65+.70*\u/360)*sin(\u)});
 \draw[very thick] (1.35,0)--(3.05,0);
 \draw[very thick,->] (1.75,0)--(2.65,0);
 \draw[very thick,domain=0:360,samples=100,variable=\u]
       plot ({4.4+(.65+.70*\u/360)*cos(180+\u)},
             {(.65+.70*\u/360)*sin(180+\u)});
 \draw[very thick,->] (3.75,0)--(4.18,0);
 \draw[blue!65!black,dashed,thick] (-.22,0)--(-.65,0);
 \draw[blue!65!black,dashed,thick,domain=0:360,samples=100,variable=\u]
       plot ({(.65+.70*\u/360)*cos(180+\u)},
             {(.65+.70*\u/360)*sin(180+\u)});
 \draw[blue!65!black,dashed,thick] (4.62,0)--(5.05,0);
 \draw[blue!65!black,dashed,thick,domain=0:360,samples=100,variable=\u]
       plot ({4.4+(.65+.70*\u/360)*cos(\u)},
             {(.65+.70*\u/360)*sin(\u)});
 \draw[blue!65!black,dashed,thick,rounded corners=12pt]
       (5.75,0)--(6.10,0)--(6.10,-1.80)--(-1.70,-1.80)
       --(-1.70,0)--(-1.35,0);
 \draw[blue!65!black,thick,->] (2.70,-1.80)--(1.70,-1.80);
 \filldraw[fill=white,draw=black,thick]
       (.38,-.08) rectangle (.54,.08);
 \filldraw[fill=white,draw=black,thick]
       (3.86,-.08) rectangle (4.02,.08);
 \node[above,fill=white,inner sep=2pt] at (2.2,.10)
       {straight gap};
 \node[below,blue!65!black] at (2.2,-1.84)
       {common complementary connector};
 \node[above] at (0,1.51) {turning disk at $s$};
 \node[above] at (4.4,1.51) {turning disk at $t$};
 \node[above,align=center] at (2.2,.78)
       {fixed endpoint\\insertions};
 \draw[->] (1.48,.77)--(.48,.12);
 \draw[->] (2.92,.77)--(3.94,.12);
\end{tikzpicture}
\caption{A standard representative, shown schematically with one full
turn at each terminal. The solid route runs from $s$ to $t$; the dashed
opposite strands and remote arc form its common complement. The small
squares mark the reference copies used by the fixed endpoint insertions.
The proof uses polygonal
spirals and inserts further reference copies before and after their
bends; these are omitted here. The turning disks have radius
$O(h(1+k))$, while the number of short pieces is $O(1+k)$.}
\label{fig:standard-representative}
\end{figure}


\subsection{The affine rule and its uniform terms}

\begin{proposition}[Endpoint potentials]\label{prop:endpoint-potentials}
For every admitted departure or arrival placement and type there are
centered state functions $\phi_s^{X,+}$ and $\phi_t^{X,-}$, allowed to
depend on $\eps$, with
\[
 \|\phi_s^{X,+}\|_{L^2(\mu_s)}\le C_{\rm short}h,
 \qquad
 \|\phi_t^{X,-}\|_{L^2(\mu_t)}\le C_{\rm short}h.
\]
They depend only on the indicated endpoint data, and not on the long
connector or its allowed small chart changes.  For each fixed connector
$e:s\to t$ of class size $k$, provided $|z_t-z_s|>Ch(1+k)$,
\begin{equation}\label{affine:transport}
 \begin{split}
 i_e^\eps(\sigma_s,\sigma_t)
   ={}&a_h^X(z_t-z_s)
       +\phi_s^{X,+}(\sigma_s)+\phi_t^{X,-}(\sigma_t)
       +b_e^\eps+\rho_e^\eps,\\
 &\|\rho_e^\eps\|_{L^1(\mu_s\otimes\mu_t)}\longrightarrow0,
 \qquad
 \limsup_{\eps\to0}|b_e^\eps|\le C_{\rm short}h(1+k).
 \end{split}
\end{equation}
Here $b_e^\eps$ is deterministic.  The constants depend on the fixed
normalized menu and short templates, but not on a thin fitted connector,
its length or a finite cover of such connectors.
\end{proposition}

\begin{proof}
If $e$ uses changed charts, first hold those charts fixed and compare it,
by Lemma~\ref{lemma:two-twists} and the common-closure rule, to the
corresponding image of its standard representative.  Apply
Lemma~\ref{lemma:chart-stability} only to this buffered representative
to return to the nominal charts.  Finally compare the nominal unsplit
representative to the inserted system of
Lemma~\ref{lemma:standard-spirals}, using its common complement.  Thus,
on the product of endpoint and auxiliary state laws,
\[
 i_e^\eps=\sum_j i_j^\eps+o_{L^1}(1).
\]
Each chain of comparisons is finite with its pattern fixed.  Narrow
bands, small positive conditioning probabilities and the length of a
finite deformation chain may enter the rate of this vanishing error;
they do not enter the short-template constant.

Average over every auxiliary state, leaving $\sigma_s,\sigma_t$ fixed.
Only the fixed first and last insertions still have a nonconstant part.
Define
\[
 \begin{split}
 \phi_s^{X,+}(\sigma_s)
   &=\EE_\mu[i_{\rm first}^\eps\mid\sigma_s]
                         -\EE_\mu i_{\rm first}^\eps,\\
 \phi_t^{X,-}(\sigma_t)
   &=\EE_\mu[i_{\rm last}^\eps\mid\sigma_t]
                         -\EE_\mu i_{\rm last}^\eps.
 \end{split}
\]
Conditional expectation followed by centering is an orthogonal
projection in $L^2$, so the fixed short-completion bounds give the stated
norms.  The definitions use the nominal fixed insertions; this is why
neither the long route nor the allowed chart changes affect them.
Conditional Jensen preserves the comparison's $o_{L^1}(1)$ error.

Let $J$ be the short pieces and $T$ the straight ones, with oriented
center displacements $d_j$.  Their total displacement is $z_t-z_s$.
For each fixed $j\in T$, Proposition~\ref{prop:affine} gives
$\EE_\mu i_j^\eps-a_h^X d_j\to0$.  Put the sum of these finitely many
deterministic deviations into $\rho_e^\eps$, together with the preceding
comparison error.  Define the remaining deterministic term by
\[
 b_e^\eps=\sum_{j\in J}
             \bigl(\EE_\mu i_j^\eps-a_h^X d_j\bigr).
\]
There are at most $C(1+k)$ such pieces.  Each has cochain norm at most
$C_{\rm short}h$, displacement $O(h)$ and the same bounded scalar
$a_h^X$.  Enlarging $C_{\rm short}$ once gives the bias bound in
\eqref{affine:transport}.  In particular a large total length of straight
legs creates no length-proportional bias: their limiting means have
already been included exactly in the affine displacement.
\end{proof}

\subsection{Geometry of the standard representatives}

It remains to construct the representative used in the endpoint formula.
We build the two terminal turning regions, verify their separation, and
then insert reference copies to isolate uniformly bounded corner pieces.

\begin{proof}[Proof of Lemma~\ref{lemma:standard-spirals}]
All lengths until the last rescaling are normalized.  Let $D_{\rm st}$
be the lower length threshold for a straight transport.  Continue the
two stubs at each terminal along opposite east and west spokes beyond
all local structures.  Set $\delta=2\pi/16$.  To make a turn $\Theta$,
including its bounded directional adjustment and integer full turns,
write
\[
 |\Theta|=N\delta+\alpha,\qquad
 N\in\Z_{\ge0},\quad 0\le\alpha<\delta.
\]
Choose a large constant $D_2$ and take $D_1=D_2$.  Starting at radius
$D_1$, use $N$ polygonal legs of angular increment $\delta$ in the sign
of $\Theta$, increasing radius by $D_2$ at each leg.  If $\alpha>0$,
append one leg of angle $\alpha$, also increasing radius by $D_2$.
The second strand is the image under rotation by $\pi$.  Both finish
with radial outward continuations.

At $s$ choose the final used spoke toward $t$.  At $t$ choose its final
used spoke, viewed outward, back toward $s$; the connecting arc traverses
this latter spiral inward.  In the integer convention of
Lemma~\ref{lemma:two-twists}, the outward constructions therefore use
\[
 \Theta_s=\theta_s+2\pi n_s,\qquad
 \Theta_t=\theta_t-2\pi n_t,
 \qquad |\theta_s|+|\theta_t|\le C.
\]
This specifies both used strands and their unused opposite companions.

We verify the geometric margins, including a last partial leg.  A leg
from radius $r$ to $r+D_2$ with angular increment
$0<\gamma\le\delta$ has length at least $D_2$, since
\[
 |(r+D_2)e^{\ii\gamma}-r|^2
   =D_2^2+2r(r+D_2)(1-\cos\gamma).
\]
Its direction relative to its initial radial ray is
\[
 \frac\gamma2+\arctan\left(
      \frac{2r+D_2}{D_2}\tan\frac\gamma2\right).
\]
This lies between $0$ and $\pi/2+\delta/2$; relative to the final radial
ray the incoming direction lies between $0$ and $\pi/2$.  Thus all
bends, including radial joins and a partial last leg, have absolute
turning angle at most
\[
 \beta_* =\pi/2+\delta/2<\pi.
\]
Reflection gives the same assertion for the opposite sign of turn.

For a full leg its radius at angle $0\le\theta\le\delta$ is
\begin{equation}\label{affine:spiral-radius}
 F_r(\theta)=
 \frac{\sin\delta}
 {\sin(\delta-\theta)/r+\sin\theta/(r+D_2)}.
\end{equation}
Uniformly for $r\ge D_2$,
\[
 \cos\delta\,r\le F_r(\theta)\le r+D_2,\qquad
 |\partial_\theta F_r(\theta)|\le C_\delta r,\qquad
 \partial_rF_r(\theta)\ge\cos(\delta/2).
\]
For the last inequality put $A=\sin(\delta-\theta)$ and
$B=\sin\theta$; Cauchy--Schwarz gives
\[
 \partial_rF_r
 =\sin\delta\frac{A/r^2+B/(r+D_2)^2}
                    {(A/r+B/(r+D_2))^2}
 \ge\frac{\sin\delta}{A+B}\ge\cos(\delta/2).
\]
In the sixteen fixed angular sectors, two distinct full legs in one
sector have starting radii differing by at least $8D_2$: a return of
one strand is offset by sixteen steps and the opposite strand by eight.
Their radial graphs are therefore separated by at least
$8D_2\cos(\delta/2)$.  For legs in adjacent sectors, the endpoint radii
at the common ray differ by at least $8D_2$, unless they are consecutive
legs of one strand.

These radial gaps give Euclidean clearance $c_0D_2$ between
nonconsecutive full legs, for a fixed $c_0>0$.  Indeed, if two points
were within $cD_2$, their radii would be comparable and their angular
separation at most $C cD_2/r$.  For small $c$, their sectors must be equal
or adjacent.  Compare their radii on a common ray, or on the intervening
step ray for adjacent sectors.  The angular derivative bound in
\eqref{affine:spiral-radius} changes the radii by only $C'cD_2$,
contradicting the preceding radial gaps.

A partial last leg need not have a uniform angular derivative bound.
Instead let $P_{r,\alpha}(\theta)$ be its radial graph.  For
$0\le\theta\le\alpha<\delta$ it lies outside the hypothetical full leg:
\[
 P_{r,\alpha}(\theta)\ge F_r(\theta).
\]
To check this, write
\[
 \frac1{P_{r,\alpha}(\theta)}
 =\frac{\cos\theta}{r}
 +\frac{\sin\theta}{r}
       \frac{r/(r+D_2)-\cos\alpha}{\sin\alpha}.
\]
For $0<a<1$, the derivative of
$(a-\cos\gamma)/\sin\gamma$ is
$(1-a\cos\gamma)/\sin^2\gamma>0$.
This proves the comparison.  Every earlier nonincident leg in the same
or an adjacent sector is on the inner side of the hypothetical full
leg, with the established gap.  The close-points argument thus uses only
the angular derivatives of that hypothetical leg and the earlier full
leg, not of the partial one.  The simultaneous partial legs of the two
strands lie in opposite sectors and are separated on the scale $D_2$.

The initial radial continuations have radius at most $D_1=D_2$.  All
turning legs after the first have radius at least
$2D_2\cos\delta>D_2$, giving a gap of order $D_2$; the first is incident
to its own initial spoke and angularly separated from the opposite one.
At the outer end, each radial continuation begins at the largest vertex
radius, whereas all earlier legs except its last incident leg have
radius at least $D_2$ less.  The opposite final leg is separated
angularly.  The central stub neighborhoods retain their fixed buffers,
and the turning legs lie beyond them when $D_2$ is sufficiently large.

The number of legs and bends is $O(1+|\Theta|)$, and their maximum radius
is $O(1+|\Theta|)$.  Thus the two terminal turning regions lie in disks
of radii $Ch(1+k)$ after rescaling.  Under the separation assumption
they are disjoint.  Join the two used final spokes by their straight
gap.  The unused spokes point away from this gap; continue them outward
and join them by a remote arc outside the two turning disks and the
used gap.  Its narrow band gives a complementary connector avoiding
all subsequent insertions.

We next fix the short corner experiments.  Let $S$ bound the reference
copy and its fixed stub structures, and choose $E_0$ large compared to
$S$.  At a corner $p$ with incoming and outgoing unit directions
$u_-,u_+$, place aligned reference copies at
$p-E_0u_-$ and $p+E_0u_+$.  The non-reversal bound gives their separation
at least $2E_0\cos(\beta_*/2)$.  Outside their stub exit squares, place
a narrow strip around the intervening two-leg path, with brief straight
continuations matching the exits.  Round the corner within a fixed
small neighborhood if desired.  The angle bound leaves positive margins
from both copies, their back stubs and adjacent straight transports.
A standalone completion joins the two unused back stubs outside a
regular neighborhood of this bounded configuration, still in a ball
of radius $C E_0$.

The bend angles range over the compact interval $[-\beta_*,\beta_*]$.
A buffered experiment at one angle persists in a neighborhood of it:
move the terminal copy together with both its stubs by a small rigid
motion, taper this motion to the identity in a disjoint collar and
deform the strip slightly.  Its fixed continuation corridors keep
their margins, so Lemma~\ref{lemma:bridge} supplies a common positive
bound in this neighborhood.  A finite subcover gives common bounds on
normalized size and density over all corner angles.  Proposition~
\ref{prop:state-cochain} now supplies one $C_{\rm short}$.  Include the
fixed first and last insertions and bounded straight-size templates in
the same finite collection of compact families.

Finally take $D_2=D_1$ much larger than $E_0,D_{\rm st}$, the central
local scales and $c_0^{-1}$.  The corner neighborhoods are disjoint.
Place copies before and after every bend as just described; after
reserving their bounded neighborhoods, all remaining aligned gaps have
length exceeding $D_{\rm st}$.  The first and last fixed insertions fit
on the initial radial parts.  Increasing the constant in the terminal
separation condition ensures the same bound for the middle gap.  Narrow
strips along all gaps give the inserted simple system.  The unchanged
remote complementary connector also closes the unsplit representative.

There are only $O(1+k)$ short pieces, each with displacement $O(h)$.
No linear bound on the total arclength is asserted or needed: the
spiral's full legs can have total length of order $(1+k)^2$ in normalized
units.  For each nominal terminal pair, type pair and pair of integers,
fix this construction once.  In the support of allowed chart changes it
uses only fixed local stubs, with no return of remote pieces.  The
auxiliary placements and straight lengths can therefore be chosen
before those changes.  Grid terminal pairs and twist pairs are
countable, the menu is finite, and the scale list is countable.
Hence the required straight lengths are countable and may be included
in the simultaneous extraction of Proposition~\ref{prop:affine}.
\end{proof}

All limits of straight means may be fixed jointly at countably many
scales for both fields.  Subsequent comparisons of individual fixed
connectors, including the finite fitted covers in
Section~\ref{sec:actual-arcs}, use only these limits and the vanishing
cochain comparison errors.  They require no additional subsequence
depending on the sampled rim.  At each fixed $h$ a finite sum of
pattern-dependent failure probabilities still vanishes as
$\eps\downarrow0$; only afterward is $h$ decreased.  The constants in
the potential and bias terms remain the common short-template constants
throughout.
